\documentclass[10pt,a4paper]{amsart}
\usepackage[margin=19mm]{geometry}
\usepackage{amsmath,amssymb,mathtools}
\usepackage[scr=rsfs,cal=euler]{mathalfa}
\usepackage{xfrac}
\usepackage[unicode,hidelinks,bookmarks=false]{hyperref}
\newcommand{\ie}{i.e.\ }
\newcommand{\cf}{cf.\ }
\newcommand{\resp}{respectively\ }
\newcommand{\Subsection}{\subsection}
\providecommand{\nobreakdash}{\nobreak\hbox{-}}
\setlength{\emergencystretch}{2em}
\multlinegap0pt
\usepackage{amssymb}
\usepackage{dsfont}
\usepackage{mathtools}
\newtheorem{thm}{Theorem}
\newcommand{\T}{\mathbb{T}}
\newcommand{\dd}{\mathop{}\!\mathrm{d}}
\DeclareMathOperator{\dv}{\dd{vol}}
\newcommand{\eps}{\varepsilon}
\newcommand{\pD}{\mathbf{D}}
\newcommand{\parenthesis}[1]{\left(#1\right)} % round bracket
\title{Periodic Solutions for a Nonlinear Dirac Equation with Soler-Type Nonlinearity}
\author{Fuping Zhang}
\date{Source: arXiv:2608.25320v1 (CC BY 4.0)}
\begin{document}
\maketitle

\noindent\emph{Source and licence.} Periodic Solutions for a Nonlinear Dirac Equation with Soler-Type Nonlinearity. Version arXiv:2608.25320v1 (CC BY 4.0), distributed by arXiv under the Creative Commons Attribution 4.0 International licence, http://creativecommons.org/licenses/by/4.0/, as recorded in the arXiv licence metadata for this exact version and shown on the abstract page https://arxiv.org/abs/2608.25320v1. Full licence notice: \texttt{licenses/arxiv-2608.25320-CC-BY-4.0.txt}.\par\smallskip

\noindent\textit{Editorial scope.} Counted unit: the article's thm thm:Soler-subthreshold (source0.tex lines 1018--1030). The article's complete original proof of it is delivered with it (source0.tex lines 1032--1329, one \texttt{proof} environment of 298 lines). No mathematics is written, repaired or re-derived: the statement and the proof are the article's own lines, in the article's own notation, and no retained line is substituted or re-flowed.  Retained uncounted as the source-local context the proof needs: lines 881--897.  Nothing was omitted from any retained span, so the package carries no omission record.  No retained line contains a citation, so the wrapper carries no bibliography and no bibliography entry.  No retained line contains a figure environment, an image-inclusion command or drawing code, so no image is delivered and the package declares no asset dependency.  Editorial formatting only, in addition: an AMS \texttt{amsart} preamble that restates the article's own declarations and macros used by the retained lines, namely the environments \texttt{thm} on separate counters, together with the standard packages those lines need.  The retained environments print in this document as Theorem 1 in order of appearance, whereas the article prints the same items with the numbering of its own full document.  The complete native source of the article as distributed in the arXiv e-print for this exact version is bundled unchanged as \texttt{sources/208567/source0.tex}.
\begin{align}\label{eq:uniform estimate for perturbation}
c_\eps(m,a)
\leq
\sup_{\substack{\psi^-\in H^{1/2,-},\,t\ge0\\
\|\psi^-\|_{L^2}<t}}
J_\eps(\psi^-+te)
&\leq
\sup_{\substack{\psi^-\in H^{1/2,-},\,t\ge0\\
\|\psi^-\|_{L^2}<t}}
J(\psi^-+te)\\
&\leq U(m,a)=
\frac{\nu-1}{\nu^{\nu^*}}
\frac{\operatorname{vol}(\mathbb T^3)}
{h_{\min}^{\frac1{\nu-1}}\,2^{\nu^*}}
(m-a)^{\frac{\nu}{\nu-1}},
\qquad\text{for every }\eps\in(0,1].
\end{align}

\begin{thm}
\label{thm:Soler-subthreshold}
Assume that
\begin{align}
U(m,a)<c_\infty(m,a)=\frac{\nu-1}{(2\nu)^{\nu^*}}\cdot\frac{1}{h_{\max}^{\frac{1}{\nu-1}}}\kappa(m,a)^{\nu^*}.
\end{align}
Let $\psi_\eps$ be a critical point of $J_\eps$ at the
minimax level $c_\eps(m,a)$. Then
\begin{align}
\sup_{\eps\in(0,1]}
\|\psi_\eps\|_{H^1(\T^3)}<+\infty.
\end{align}
\end{thm}

\begin{proof}
We first prove that there exists $C=C(\mathbb{T}^3,m,a,\nu)>0$, independent of $\eps\in(0,1]$, such that
\begin{align}
\|\psi_\eps\|_{L^3(\T^3)}\leq C.
\end{align}
Suppose for contradiction that the assertion fails. Then there exists a sequence $\eps_n\to0^+$ with
\begin{align}
\lambda_n :=
\|\psi_{\eps_n}\|_{L^3}
\to+\infty.
\end{align}
Set
\begin{align}
\varphi_n=\frac{\psi_{\eps_n}}{\lambda_n},
\end{align}
so that
\begin{align}
\|\varphi_n\|_{L^3}=1.
\end{align}


The normalized equations
\begin{align}\label{eq:normalized eqn}
    \pD\varphi_n-a\varphi_n=2\nu h(x)\lambda_n^{2\nu-2}|\bar{\varphi_n}\varphi_n|^{\nu-2}(\bar{\varphi_n}\varphi_n)\gamma^0\varphi_n + 2\nu\eps_n |\psi_{\eps_n}|^{2\nu -2 }\varphi_n, \qquad \mbox{ on } \; \T^3.
\end{align}
Define
\begin{align}
V_n
:=
2\nu h(x)\lambda_n^{2\nu-2}
\left|\bar{\varphi_n}\varphi_n\right|^{\nu-2}\bar{\varphi_n}\varphi_n.
\end{align}
Then
\begin{align}
(\pD-a)\varphi_n
=
V_n
\gamma^0\varphi_n
+
R_n,
\end{align}
where
\begin{align}
R_n
:=
2\nu\eps_n
|\psi_{\eps_n}|^{2\nu-2}\varphi_n.
\end{align}
We have
\begin{align}
   \|R_n\|_{L^{\frac{2\nu}{2\nu -1}}}=\|2\nu\eps_n |\psi_{\eps_n}|^{2\nu -2 }\varphi_n \|_{L^{\frac{2\nu}{2\nu -1}}}
    = & \frac{2\nu \eps_n^{\frac{1}{2\nu}} }{\lambda_n}
     \parenthesis{ \int_{\T^3} \left|\eps_n^{\frac{2\nu-1}{2\nu}} |\psi_{\eps_n}|^{2\nu -1} \right|^{\frac{2\nu}{2\nu-1}} \dv   }^{\frac{2\nu-1}{2\nu}}\\
    \leq& \frac{2\nu \eps_n^{\frac{1}{2\nu}} }{\lambda_n}
    \parenthesis{ \int_{\T^3} \eps_n |\psi_{\eps_n}|^{2\nu}\dv}^{\frac{2\nu-1}{2\nu}} \\
    \leq& \frac{2\nu \eps_n^{\frac{1}{2\nu}} }{\lambda_n}\parenthesis{ \frac{U(m,a)}{\nu-1} }^{\frac{2\nu-1}{2\nu}}
\end{align}
thanks to~\eqref{eq:uniform estimate for perturbation}.
 Thus, $R_n\to0
\quad\text{in }L^{\frac{2\nu}{2\nu-1}}$.




\begin{align}
\|V_n\|^{\nu^*}_{L^{\nu^*}}&=(2\nu)^{\nu^*}\lambda_n^{(2\nu-2)\nu^*}\int_{\T^3}h(x)^{\nu^*}|\bar{\varphi_n}\varphi_n|^{(\nu-1)\nu^*}\dv\\
&=(2\nu)^{\nu^*}\lambda_n^{2\nu}\int_{\T^3}h(x)^{\nu^*}|\bar{\varphi_n}\varphi_n|^{\nu}\dv\\
&\leq (2\nu)^{\nu^*}h_{\max}^{\nu^*-1}\int_{\T^3}h(x)|\bar{\psi_{\eps_n}}\psi_{\eps_n}|^{\nu}\dv\\
&\leq \frac{(2\nu)^{\nu^*}h_{\max}^{\nu^*-1}}{\nu-1}\cdot U(m,a).
\end{align}
Thus ${V_n}$ is bounded in $L^{\nu^*}$. By weak compactness, passing to a subsequence,
\begin{align}
V_n\rightharpoonup V
\quad\text{weakly in }L^{\nu^*}.
\end{align}

By H\"{o}lder's inequality, using $\frac1{\nu'}=\frac1{\nu^*}+\frac13$, $\|V_n\|_{\nu^*}\le C$ and $\|\varphi_n\|_{L^3}=1$, since $ \nu'>\frac{2\nu}{2\nu-1}>\frac{3}{2}$ and $\T^3$ has finite volume, we obtain \begin{align}
\|V_n\varphi_n\|_{L^{ \nu'}}\leq C \Rightarrow\|V_n\varphi_n\|_{L^{\frac{2\nu}{2\nu-1}}}\leq C.
\end{align}
By elliptic regularity,
\begin{align}
\|\varphi_n\|_{W^{1,\frac{2\nu}{2\nu-1}}}
&\le
C\left(
\|V_n\varphi_n\|_{L^{\frac{2\nu}{2\nu-1}}}
+\|R_n\|_{L^{\frac{2\nu}{2\nu-1}}}
\right)\\
&\le C.\qquad \frac{2\nu}{2\nu-1}>\frac{3}{2}.
\end{align}
Therefore, after extracting a further subsequence,
\begin{align}
\varphi_n\rightharpoonup\varphi
\quad\text{in }W^{1,\frac{2\nu}{2\nu-1}}.
\end{align}
Since $1<\nu<\frac{3}{2}$, the embedding~$W^{1,\frac{2\nu}{2\nu-1}}(\T^3)\Subset L^{3}(\T^3)$ is compact. Therefore,
\begin{align}
\varphi_n\to\varphi
\quad\text{strongly in }L^3.
\end{align}
We know $\|\varphi\|_{L^3}=1$, so
\begin{align}
\varphi\ne0.
\end{align}
We have 
\begin{align}
\bar\varphi_n\varphi_n=
\lambda_n^{-2}
\bar\psi_{\eps_n}\psi_{\eps_n}.
\end{align}
We now prove the null condition. We have
\begin{align}\label{the null condition}
\int_{\T^3}|\bar{\varphi_n}\varphi_n|^{\nu}\dv&=\frac{\int_{\T^3}|\bar{\psi_{\eps_n}}\psi_{\eps_n}|^{\nu}\dv}{\lambda_n^{2\nu}}\\
&\leq \frac{C}{h_{\min}\lambda_n^{2\nu}}\to0.
\end{align}
Since $2\nu<3$ and $\T^3$ has finite volume, strong convergence in $L^3$ implies strong convergence in $L^{2\nu}$. Therefore,
\begin{align}
\bar\varphi_n\varphi_n\to\bar\varphi\varphi
\quad\text{in }L^\nu.
\end{align}
Combined with \eqref{the null condition}, we know
\begin{align}
\bar\varphi\varphi=0
\quad\text{a.e. on }\T^3.
\end{align}
For any smooth test spinor $\eta$, we have:
\begin{align}
\int_{\T^3}
\langle(\pD-a)\varphi_n,\eta\rangle\dv=
\int_{\T^3}
\langle V_n\gamma^0\varphi_n,\eta\rangle\dv
+
\int_{\T^3}\langle R_n,\eta\rangle\dv.
\end{align}
The left-hand side converges to
\begin{align}
\int_{\T^3}\langle(\pD-a)\varphi,\eta\rangle\dv.
\end{align}
Because $\varphi_n\to\varphi$ strongly in $L^3$ and $V_n\rightharpoonup V$ in $L^{\nu^*}$,
\begin{align}
\int_{\T^3}
\langle V_n\gamma^0\varphi_n-V\gamma^0\varphi,\eta\rangle\dv&=\int_{\T^3}
\langle V_n\gamma^0(\varphi_n-\varphi),\eta\rangle\dv+\int_{\T^3}
\langle (V_n-V)\gamma^0\varphi,\eta\rangle\dv.
\end{align}
The first term tends to zero because
\begin{align}
\|V_n\|_{L^{\nu^*}}\le C,\qquad
\|\varphi_n-\varphi\|_{L^3}\to0.
\end{align}
For the second term,
$\langle\gamma^0\varphi,\eta\rangle\in L^\nu$,
since $L^3\subset L^\nu$ on the finite-volume torus, and therefore weak convergence
$V_n\rightharpoonup V\quad\text{in }L^{\nu^*}$
makes it tend to zero. Thus,
\begin{align}
V_n\gamma^0\varphi_n
\rightharpoonup
V\gamma^0\varphi\quad \text{in the distributional sense}.
\end{align}
Hence
\begin{align}
(\pD-a)\varphi=V\gamma^0\varphi
\end{align}
holds in the weak sense. We can recover the regularity required in the definition of
\begin{align}
\kappa:
V\in L^{\nu^*},\qquad \varphi\in L^3
\quad\Longrightarrow\quad
V\varphi\in L^{\nu'},
\end{align}
and elliptic regularity gives
\begin{align}
\varphi\in W^{1,\nu'}.
\end{align}
Thus, the pair $(\varphi,V)$ is admissible in the definition of $\kappa(m,a)$. By the definition of $\kappa(m,a)$ and weak lower semicontinuity of the norm,
\begin{align}
\kappa(m,a)^{\nu^*}
&\leq
\|V\|_{L^{\nu^*}}^{\nu^*}\\
&\leq
\liminf_{n\to\infty}
\|V_n\|_{L^{\nu^*}}^{\nu^*}\\
&\leq
(2\nu)^{\nu^*}
h_{\max}^{\nu^*-1}
\liminf_{n\to\infty}
\int_{\T^3}
h(x)|\bar\psi_{\eps_n}\psi_{\eps_n}|^\nu\dv\\
&\leq \frac{(2\nu)^{\nu^*}
h_{\max}^{\nu^*-1}}{\nu-1}\liminf_{n\to\infty}c_{\eps_n}.
\end{align}
It follows that
\begin{align}
\liminf_{n\to\infty}c_{\eps_n}(m,a)
&\geq
\frac{\nu-1}{(2\nu)^{\nu^*}}
\frac{1}{h_{\max}^{\frac{1}{\nu-1}}}
\kappa(m,a)^{\nu^*}\\
&=:c_\infty(m,a).
\end{align}



We have
\begin{align}
J_\eps(\psi)=J_0(\psi)-\eps\|\psi\|_{L^{2\nu}}^{2\nu}
\le J_0(\psi).
\end{align}
By the upper-bound estimate established above,
$\sup_{\mathscr C_1(R)}J_0\le U(m,a)$; moreover $J_\eps\le J_0$. For each $\eps>0$, the identity map belongs to the admissible class $\Gamma_\eps$, hence
\begin{align}
c_\eps(m,a)
\le
\sup_{\mathscr C_1(R)}J_\eps
\le U(m,a).
\end{align}
By the hypothesis of the theorem,
\begin{align}
U(m,a)<c_\infty(m,a),
\end{align}
we obtain
\begin{align}
c_\infty(m,a)
\le
\liminf_{n\to\infty}c_{\eps_n}(m,a)
\le
\limsup_{n\to\infty}c_{\eps_n}(m,a)
\le U(m,a)
<c_\infty(m,a),
\end{align}
a contradiction.

Therefore
\begin{align}
\sup_{\eps\in(0,1]}
\|\psi_\eps\|_{L^3}<+\infty.
\end{align}

By the uniform $L^3$-estimate, both of the nonlinear terms in the equation
\begin{align}
(\pD-a)\psi_\eps
=
2\nu h(x)|\overline{\psi_\eps}\psi_\eps|^{\nu-2}
(\overline{\psi_\eps}\psi_\eps)\gamma^0\psi_\eps
+
2\nu\eps|\psi_\eps|^{2\nu-2}\psi_\eps,
\end{align}
are bounded pointwise by $C|\psi_\eps|^{2\nu-1}$. Hence
\begin{align}
\|(\pD-a)\psi_\eps\|_{L^{\frac{3}{2\nu-1}}}
\le C.
\end{align}
Since $1<\nu<3/2$, we have
\begin{align}
\frac32<\frac{3}{2\nu-1}<3.
\end{align}
Elliptic regularity and Sobolev embedding therefore yield
\begin{align}
\|\psi_\eps\|_{W^{1,\frac{3}{2\nu-1}}}\le C,
\qquad
\|\psi_\eps\|_{L^{\frac{3}{2\nu-2}}}\le C.
\end{align}

More generally, suppose that $\psi_\eps$ is bounded in $L^r$ for
some $r\ge3$. Then
\begin{align}
(\pD-a)\psi_\eps
\end{align}
is bounded in $L^{r/(2\nu-1)}$. As long as
$r/(2\nu-1)<3$, elliptic regularity followed by Sobolev embedding gives a
bound in $L^{r_{\mathrm{new}}}$, where
\begin{align}
\frac1{r_{\mathrm{new}}}
=
\frac{2\nu-1}{r}-\frac13.
\end{align}
Since $\nu<3/2$, this iteration strictly increases the integrability
exponent. Consequently, after finitely many steps, we reach an exponent
$r$ satisfying
\begin{align}
r\ge 2(2\nu-1).
\end{align}
At this stage,
\begin{align}
\frac{r}{2\nu-1}\ge2.
\end{align}
Since $\T^3$ has finite volume, $L^\frac{r}{2\nu-1}(\T^3)\hookrightarrow L^2(\T^3)$, hence
\begin{align}
\|(\pD-a)\psi_\eps\|_{L^2}\le C.
\end{align}
If at some stage $r/(2\nu-1)\ge3$, then $r\ge3(2\nu-1)>2(2\nu-1)$, and the iteration stops. The $L^2$-boundedness of $\psi_\eps$ and the elliptic estimate for
$\pD-a$ then imply
\begin{align}
\|\psi_\eps\|_{H^1}\le C,
\end{align}
with $C$ independent of $\eps\in(0,1]$.

\end{proof}

\end{document}
